\subsection{Surface Integrals}

\content{

}{% presentation
\vspace{4in}
}{% nopresentation
}{% comments
Go over the setup for the surface integral of a scalar valued function formula
ends up being 
$$ \sum\sum f(P_{ij}^*)\Delta S_{ij}. $$
so the integral is 
$$ \iint_S f(x,y,z) \, dS . $$
}

\content{
\begin{definition} The surface integral of $f(x,y,z)$ over $S$ is given by 
$$ \iint_S f(x,y,z)\, dS = \iint_D f(x,y,z)|\vec{r}_u\times \vec{r}_v|\,dA. $$
\end{definition}
}{% presentation
}{% nopresentation
}{% comments
}

\content{
\begin{example} Evaluate $\iint_S x^2\,dS$, where $S$ is the unit sphere. 
\end{example}
}{% presentation
\vspace{3in}
}{% nopresentation
}{% comments
As previously computed
$\vec{r}_\theta\times\vec{r}_\phi =\la \sin^2\phi\cos\theta,
\sin^2\phi\sin\theta, \sin\phi\cos\phi\ra. $
}


\content{
\begin{example} Evaluate $\iint_S z\,dS$, where $S$ is the surface whose sides
$S_1$ are given by the cylinder $x^2+y^2=1$, whose bottom $S_2$ is the disk
$x^2+y^2\leq 1$ in the plane $z=0$, and whose top $S_3$ is the part of the plane
$z=1+x$ that lies above $S_2$.
\end{example}
}{% presentation
\vspace{4in}
}{% nopresentation
}{% comments
$\iint_{S_1} z = 3\pi/2$, $\iint_{S_2}z = 0$, and $\iint_{S_3} z =
(3/2+\sqrt{2})\pi$.
}

\subsubsection{Oriented Surfaces}

\content{
\begin{definition} An oriented surface is a surface on which it is possible to
choose a unit normal vector $\vec{n}$ at every point so that $\vec{n}$ varies
continuously over $S$. 
\end{definition}
}{% presentation
\vspace{3in}
}{% nopresentation
}{% comments
Mention that all of the surfaces we look at in this course will be orientable,
but nonorientable surfaces do exist (Mobius strip)
}

\subsubsection{Surface Integrals of Vector Fields and Flux}

\content{

}{% presentation
\vspace{4in}
}{% nopresentation
}{% comments
Go over how the flux integral is derived the amount of flow through the surface
$S$ in the direction of $\vec{n}$ (normal vector) in a small $(u,v)$ rectangle
on the surface is approximately equal to 
$$ \vec{F}\cdot \vec{n} A(S_{ij}) $$
Leading us to 
$$ \iint_S \vec{F}\cdot \vec{n}\, dS$$
}

\content{
\begin{definition} If $\vec{F}$ is a continuous vector field defined on an
oriented surface $S$ with unit normal vector $\vec{n}$, then the usrface
integral of $\vec{F}$ over $S$ is 
$$ \iint_S \vec{F}\cdot d\vec{S} = \iint_S \vec{F}\cdot \vec{n}\,dS. $$
\end{definition}
}{% presentation
}{% nopresentation
}{% comments
Go over how this simplifies to 
$$ \iint_D \vec{F}\cdot (\vec{r}_u\times \vec{r}_v)\,dA $$ 
if the surface is parametrized by $\vec{r}(u,v)$.
}

\content{
\begin{example} Find the flux of the vector field $\vec{F}(x,y,z) = \la z, y,
x\ra$ across the unit sphere $x^2+y^2+z^2=1$.
\end{example}
}{% presentation
\vspace{4in}
}{% nopresentation
}{% comments
As previously computed, for this example 
$\vec{r}_\theta\times\vec{r}_\phi =\la \sin^2\phi\cos\theta,
\sin^2\phi\sin\theta, \sin\phi\cos\phi\ra. $

Mention that in some cases, the surface may have to be treated in parts, just as
a previous example of a scalar surface integral. 
}
